
\appendix
\section{Phụ lục}
\label{appendix:overall}

\subsection{Lý do lựa chọn kiến trúc và hồ sơ quyết định kiến trúc}\label{appendix:architecture-rationale}

Các quyết định kiến trúc của IRMS được đặt trên ba cơ sở: bối cảnh vận hành của nhà hàng, yêu cầu chức năng theo từng miền nghiệp vụ và yêu cầu phi chức năng ở mức toàn dự án. Kiến trúc chủ đạo là \textbf{service-based architecture theo miền nghiệp vụ}; \textbf{event-driven processing} đóng vai trò bổ trợ cho các hệ quả nền như thông báo, cập nhật mô hình đọc báo cáo, kiểm toán và cảnh báo tồn kho. Cách lựa chọn này giữ tuyến giao dịch chính ngắn, đồng thời vẫn cho phép các hoạt động hỗ trợ được xử lý linh hoạt phía sau.

\subsubsection{So sánh các kiểu kiến trúc ứng viên}

{\small
\setlength{\LTleft}{0pt}
\setlength{\LTright}{0pt}
\renewcommand{\arraystretch}{1.15}
\begin{longtable}{>{\raggedright\arraybackslash}p{3.1cm} >{\raggedright\arraybackslash}p{4.5cm} >{\raggedright\arraybackslash}p{6.4cm}}
\caption{So sánh các kiểu kiến trúc ứng viên cho IRMS}\label{tab:appendix-architecture-candidates}\\
\toprule
\textbf{Kiểu kiến trúc} & \textbf{Điểm phù hợp} & \textbf{Giới hạn trong bối cảnh IRMS} \\
\midrule
\endfirsthead
\toprule
\textbf{Kiểu kiến trúc} & \textbf{Điểm phù hợp} & \textbf{Giới hạn trong bối cảnh IRMS} \\
\midrule
\endhead
Kiến trúc phân lớp tập trung thuần & Dễ khởi tạo và đơn giản cho các thao tác nội bộ. & Dễ gom các nghiệp vụ khác bản chất vào cùng một tầng kỹ thuật, làm mờ ranh giới giữa đặt bàn, gọi món, bếp, thanh toán, tồn kho và báo cáo. \\
\addlinespace
Service-based theo miền nghiệp vụ & Phù hợp với cách IRMS chia trách nhiệm theo các năng lực vận hành của nhà hàng; mỗi service có ranh giới và hợp đồng giao tiếp rõ hơn. & Cần quản lý chặt hợp đồng giữa các service và tránh phụ thuộc vòng giữa các miền nghiệp vụ. \\
\addlinespace
Microservices đầy đủ & Cô lập mạnh về triển khai và hỗ trợ vòng đời phát hành riêng. & Chi phí vận hành, giám sát, quản trị giao dịch phân tán và quản lý hợp đồng dịch vụ cao hơn nhu cầu của phạm vi hiện tại. \\
\addlinespace
Event-driven thuần & Phù hợp cho phát tán sự kiện, xử lý nền và cập nhật mô hình đọc. & Không phù hợp cho toàn bộ hệ thống vì tạo đơn, lập hóa đơn, thanh toán và hoàn tiền cần phản hồi tức thời và trạng thái lỗi rõ ràng. \\
\addlinespace
Service-based kết hợp event-driven & Giữ ranh giới nghiệp vụ bằng service-based architecture, đồng thời dùng event-driven cho các hệ quả nền như thông báo, báo cáo, kiểm toán và cảnh báo tồn kho. & Cần phân biệt rõ thao tác nào đồng bộ, thao tác nào bất đồng bộ; các luồng nền cần chống xử lý lặp và có khả năng theo dõi trạng thái. \\
\bottomrule
\end{longtable}
\normalsize
}

Từ các phương án trên, service-based architecture là cấu trúc chính vì bám sát các miền trách nhiệm đã nêu trong Chương~1. Event-driven processing không thay thế cấu trúc này mà bổ sung một cơ chế phối hợp cho các hệ quả không cần chặn thao tác của nhân viên. Sự kết hợp này tạo ra một kiến trúc vừa rõ ràng ở tuyến giao dịch chính, vừa đủ linh hoạt cho thông báo, báo cáo, kiểm toán và cảnh báo tồn kho.

\subsubsection{Hồ sơ quyết định kiến trúc}

\begin{adrbox}[title={ADR-0001: Kiến trúc chủ đạo}]
\begin{description}[leftmargin=3.0cm,style=nextline]
  \item[Quyết định] Chọn \textbf{service-based architecture theo miền nghiệp vụ} làm kiến trúc chính; dùng \textbf{event-driven processing} cho các tác vụ nền và hệ quả phụ.
  \item[Cơ sở lựa chọn] Gọi món, điều phối bếp và thanh toán cần phản hồi nhanh, trong khi báo cáo, thông báo, kiểm toán và cảnh báo tồn kho có thể xử lý sau giao dịch chính.
  \item[Hệ quả] Kiến trúc phải phân biệt rõ luồng đồng bộ và bất đồng bộ; các luồng nền cần chống xử lý lặp, hỗ trợ thử lại và có mã liên kết để truy vết.
  \item[Không chọn] Kiến trúc phân lớp tập trung làm mờ ranh giới nghiệp vụ; microservices đầy đủ vượt quá mức cần thiết; event-driven thuần không phù hợp với thao tác cần phản hồi tức thời.
\end{description}
\end{adrbox}

\begin{adrbox}[title={ADR-0002: Điểm vào thống nhất}]
\begin{description}[leftmargin=3.0cm,style=nextline]
  \item[Quyết định] Dùng \textsf{Frontend Node Server} để phục vụ giao diện React và chuyển tiếp \textsf{/api}; backend dùng \textsf{API Gateway} làm điểm vào chung cho các service.
  \item[Cơ sở lựa chọn] IRMS có nhiều màn hình vận hành như gọi món, bếp, thu ngân và quản lý. Các giao diện này cần một điểm vào ổn định thay vì phụ thuộc vào vị trí của từng service nội bộ.
  \item[Hệ quả] Giao tiếp ở biên hệ thống được chuẩn hóa; thay đổi bên trong backend ít ảnh hưởng đến giao diện người dùng.
\end{description}
\end{adrbox}

\begin{adrbox}[title={ADR-0003: REST, RabbitMQ và Outbox}]
\begin{description}[leftmargin=3.0cm,style=nextline]
  \item[Quyết định] Dùng \textsf{REST} đồng bộ cho thao tác tuyến đầu; dùng \textsf{RabbitMQ}, \textsf{outbox\_events}, \textsf{processed\_events} và bộ xử lý nền cho các tác vụ sau giao dịch.
  \item[Cơ sở lựa chọn] Tạo đơn, lập hóa đơn và thanh toán cần phản hồi tức thời. Gửi thông báo, cập nhật báo cáo, ghi kiểm toán và xử lý lại khi lỗi không nên làm dài tuyến giao dịch chính.
  \item[Hệ quả] Cần kiểm soát thời điểm phát sự kiện sau giao dịch, tránh xử lý lặp và quan sát được trạng thái của hàng đợi.
  \item[Ví dụ] Tuyến đồng bộ gồm tạo đơn, lập hóa đơn và thanh toán. Tuyến bất đồng bộ gồm \textsf{OrderConfirmed}, \textsf{LowStockDetected} và yêu cầu làm mới báo cáo.
\end{description}
\end{adrbox}

\begin{adrbox}[title={ADR-0004: PostgreSQL và mô hình đọc}]
\begin{description}[leftmargin=3.0cm,style=nextline]
  \item[Quyết định] Dùng một \textsf{PostgreSQL} vật lý duy nhất trong runtime hiện tại; bên trong phân tách theo nhóm bảng logic cho dữ liệu giao dịch, outbox/inbox, định danh, kiểm toán và mô hình đọc báo cáo.
  \item[Cơ sở lựa chọn] Các giao dịch nghiệp vụ cần nhất quán mạnh, trong khi báo cáo cần tối ưu truy vấn đọc và có thể chấp nhận độ trễ có kiểm soát.
  \item[Hệ quả] Không vẽ hoặc mô tả nhiều database vật lý. Mô hình đọc là nhóm bảng logic trong cùng \textsf{PostgreSQL}, được cập nhật qua sự kiện hoặc tiến trình nền.
\end{description}
\end{adrbox}

\begin{adrbox}[title={ADR-0005: Adapter biên nội bộ}]
\begin{description}[leftmargin=3.0cm,style=nextline]
  \item[Quyết định] Cô lập thanh toán, thông báo và phát hành biên lai sau các adapter nội bộ như \textsf{PaymentGateway}, \textsf{NotificationChannelAdapter} và \textsf{ReceiptDeliveryAdapter}.
  \item[Cơ sở lựa chọn] Các kênh thanh toán, thông báo và biên lai dễ thay đổi theo chính sách vận hành. Lõi nghiệp vụ cần giữ ổn định khi thay đổi kênh hoặc thiết bị biên.
  \item[Hệ quả] Thanh toán, thông báo và phát hành biên lai được đặt sau hợp đồng trừu tượng; tài liệu phân biệt rõ adapter nội bộ với hệ thống ngoài độc lập.
\end{description}
\end{adrbox}

\begin{adrbox}[title={ADR-0006: Runtime Docker Compose}]
\begin{description}[leftmargin=3.0cm,style=nextline]
  \item[Quyết định] Dùng \textsf{Docker Compose} làm mô hình triển khai hiện trạng cho frontend, api-gateway, các service \textsf{Spring Boot}, \textsf{PostgreSQL}, \textsf{RabbitMQ} và migration container.
  \item[Cơ sở lựa chọn] Phạm vi hiện tại cần một runtime cục bộ, rõ thành phần và đủ để thể hiện cách hệ thống vận hành theo service. Các hạ tầng sẵn sàng cao được xem là hướng mở rộng, không phải phần triển khai hiện trạng.
  \item[Hệ quả] Góc nhìn triển khai bám vào các container hiện có; không mô tả Kubernetes, cân bằng tải hoặc dự phòng nóng như năng lực đã triển khai.
\end{description}
\end{adrbox}

\begin{adrbox}[title={ADR-0007: Phân quyền và kiểm toán}]
\begin{description}[leftmargin=3.0cm,style=nextline]
  \item[Quyết định] Áp dụng phân quyền theo vai trò tại lớp định danh và lớp biên; kết hợp \textsf{AuthorizationPolicy} trong miền nghiệp vụ; ghi \textsf{AuditLog} cho hành động nhạy cảm.
  \item[Cơ sở lựa chọn] Hoàn tiền, thay đổi giá, điều chỉnh tồn kho, truy cập nhật ký kiểm toán và phân quyền người dùng cần kiểm soát chặt chẽ.
  \item[Hệ quả] Kiểm soát truy cập không chỉ nằm ở lớp kỹ thuật biên mà còn xuất hiện trong quy tắc nghiệp vụ; nhật ký kiểm toán cần bảo đảm khả năng truy vết.
\end{description}
\end{adrbox}

\subsection{Liên kết giữa phụ lục và các phần chính}\label{appendix:traceability}

Phụ lục này chỉ đóng vai trò giải thích các quyết định đã được dùng trong thân tài liệu. Các nội dung dưới đây tránh lặp lại toàn bộ phần chính, mà chỉ chỉ ra nơi từng quyết định được thể hiện rõ nhất.

\begin{table}[H]
\centering
\caption{Liên kết giữa các phần chính và phụ lục biện minh}
\label{tab:appendix-traceability}
\small
\begin{tabularx}{\textwidth}{|p{3.7cm}|p{4.1cm}|X|}
\hline
\textbf{Phần chính} & \textbf{Điểm cần giải thích} & \textbf{Phụ lục liên quan} \\
\hline
Chương~1 & Kiến trúc chủ đạo xuất phát từ yêu cầu chức năng và phi chức năng. & A.1 giải thích vì sao service-based là kiến trúc chính và event-driven là cơ chế bổ trợ. \\
\hline
Mục~3.1 & Sơ đồ tổng quan dùng service-based, event-driven, một PostgreSQL và các adapter biên. & A.1 và A.4 liên kết các quyết định chính với kiến trúc tổng quan. \\
\hline
Mục~3.2--3.7 & Ba góc nhìn chính gồm Module View, Component-and-Connector View và Allocation/Deployment View. & A.3 giải thích vai trò từng góc nhìn và lý do không xem bản đồ năng lực logic là một view chính độc lập. \\
\hline
Chương~4 & Sơ đồ lớp và sơ đồ hoạt động cụ thể hóa các ranh giới service, adapter và luồng xử lý. & A.4 giải thích các thành phần nền như \textsf{PaymentGateway}, \textsf{NotificationChannelAdapter}, \textsf{ReceiptDeliveryAdapter}, phân quyền và kiểm toán. \\
\hline
\end{tabularx}
\normalsize
\end{table}

\subsection{Lý do chọn các góc nhìn và loại sơ đồ}\label{appendix:view-rationale}

Ba góc nhìn chính được chọn để tránh trộn lẫn cấu trúc tĩnh, cấu trúc runtime và cấu trúc triển khai. Bản đồ năng lực logic được giữ như nội dung hỗ trợ đọc nghiệp vụ, không phải một góc nhìn kiến trúc chính độc lập.

\begin{table}[H]
\centering
\caption{Vai trò của các góc nhìn và loại sơ đồ trong IRMS}
\label{tab:appendix-view-rationale}
\small
\begin{tabularx}{\textwidth}{|p{3.5cm}|p{3.1cm}|X|}
\hline
\textbf{Loại sơ đồ/góc nhìn} & \textbf{Vị trí chính} & \textbf{Vai trò} \\
\hline
Module View & Mục~3.5 & Mô tả các đơn vị hiện thực chính, trách nhiệm module, quan hệ phụ thuộc tĩnh và mapping giữa năng lực nghiệp vụ với cấu trúc backend. \\
\hline
Component-and-Connector View & Mục~3.6 & Mô tả thành phần runtime, port, interface và connector giữa frontend, gateway, service nghiệp vụ, PostgreSQL, RabbitMQ và các adapter biên. \\
\hline
Allocation/Deployment View & Mục~3.7 & Mô tả cách các thành phần phần mềm được phân bổ lên container, node runtime, cơ sở dữ liệu, broker, volume và cấu hình triển khai. \\
\hline
Bản đồ năng lực logic & Mục~3.4 & Hỗ trợ chuyển từ yêu cầu nghiệp vụ sang ranh giới module; không thay thế ba góc nhìn kiến trúc chính. \\
\hline
Class Diagram & Mục~4.1--4.2 & Làm rõ cấu trúc lớp, interface, service, policy, strategy và adapter ở mức thiết kế chi tiết. \\
\hline
Activity Diagram & Mục~4.3 & Kiểm tra luồng xử lý chi tiết của từng nhóm use case, đặc biệt các điểm đồng bộ, bất đồng bộ, kiểm quyền và ghi nhận kiểm toán. \\
\hline
\end{tabularx}
\normalsize
\end{table}

\subsection{Lý do chọn các thành phần và loại hình kiến trúc nền}\label{appendix:component-rationale}

Các thành phần dưới đây nối yêu cầu vận hành ở Chương~1 với kiến trúc tổng quan và thiết kế chi tiết. Nội dung được trình bày theo vai trò kiến trúc, tránh mô tả như danh sách file triển khai.

{\small
\setlength{\LTleft}{0pt}
\setlength{\LTright}{0pt}
\renewcommand{\arraystretch}{1.16}
\begin{longtable}{>{\raggedright\arraybackslash}p{3.25cm} >{\raggedright\arraybackslash}p{5.4cm} >{\raggedright\arraybackslash}p{5.2cm}}
\caption{Lý do chọn các thành phần kiến trúc nền của IRMS}\label{tab:appendix-component-rationale-core}\\
\toprule
\textbf{Thành phần / loại hình} & \textbf{Lý do dùng trong IRMS} & \textbf{Liên hệ trong kiến trúc} \\
\midrule
\endfirsthead
\toprule
\textbf{Thành phần / loại hình} & \textbf{Lý do dùng trong IRMS} & \textbf{Liên hệ trong kiến trúc} \\
\midrule
\endhead
\textsf{ReactJS} cho lớp giao diện & Phù hợp với các màn hình thao tác theo vai trò như lễ tân, phục vụ, bếp, thu ngân và quản lý; hỗ trợ cập nhật trạng thái giao diện và điều hướng nhất quán. & Xuất hiện ở biên người dùng trong kiến trúc tổng quan và các kịch bản nghiệp vụ. \\
\addlinespace
\textsf{Frontend Node Server} và \textsf{api-gateway} & Tạo lớp biên thống nhất: \textsf{Frontend Node Server} phục vụ React và chuyển tiếp \textsf{/api}, còn \textsf{api-gateway} định tuyến yêu cầu đến các service nội bộ. & Gắn với quyết định điểm vào thống nhất và các sơ đồ tổng quan, thành phần và triển khai. \\
\addlinespace
Service-based architecture & Phân rã hệ thống theo các miền nghiệp vụ như đặt bàn, gọi món, bếp, thanh toán, tồn kho, báo cáo và kiểm toán. & Là lựa chọn kiến trúc chính; được phản ánh trong Module View, Component-and-Connector View và sơ đồ tổng quan. \\
\addlinespace
Event-driven processing & Tách các hệ quả không cần phản hồi ngay khỏi tuyến giao dịch nóng, ví dụ thông báo, kiểm toán, cập nhật mô hình đọc và cảnh báo tồn kho. & Là cơ chế bổ trợ cho service-based architecture; xuất hiện trong C\&C View, các bộ xử lý nền và luồng sự kiện. \\
\addlinespace
\textsf{RabbitMQ}, outbox và bộ xử lý nền & Bảo đảm sự kiện được phát sau khi giao dịch chính đã ghi nhận; hỗ trợ thử lại và xử lý nền mà không làm nghẽn thao tác gọi món, bếp hoặc thanh toán. & Gắn với quyết định kiểu kết nối và các sơ đồ thành phần, kết nối, triển khai. \\
\addlinespace
\textsf{PostgreSQL} vật lý duy nhất & Bảo đảm nhất quán dữ liệu mạnh cho xác nhận đơn, thanh toán, hoàn tiền và kiểm toán. & Các nhóm dữ liệu như outbox, inbox, mô hình đọc báo cáo, định danh và kiểm toán là nhóm bảng logic trong cùng \textsf{PostgreSQL}, không phải database riêng. \\
\addlinespace
Mô hình đọc cho báo cáo & Tách luồng đọc báo cáo khỏi dữ liệu giao dịch nóng, giúp truy vấn dashboard không ảnh hưởng trực tiếp đến gọi món, hóa đơn hoặc thanh toán. & Được cập nhật thông qua sự kiện hoặc tiến trình nền và được mô tả như một nhóm bảng logic trong \textsf{PostgreSQL}. \\
\addlinespace
\textsf{Docker Compose} runtime & Phù hợp với phạm vi triển khai cục bộ gồm frontend, api-gateway, các service \textsf{Spring Boot}, \textsf{PostgreSQL}, \textsf{RabbitMQ} và migration container. & Gắn với Deployment/Allocation View; các hạ tầng sẵn sàng cao được đặt ở định hướng mở rộng thay vì mô hình hiện trạng. \\
\bottomrule
\end{longtable}
\normalsize
}

{\small
\setlength{\LTleft}{0pt}
\setlength{\LTright}{0pt}
\renewcommand{\arraystretch}{1.16}
\begin{longtable}{>{\raggedright\arraybackslash}p{3.25cm} >{\raggedright\arraybackslash}p{5.4cm} >{\raggedright\arraybackslash}p{5.2cm}}
\caption{Lý do chọn các adapter biên và cơ chế kiểm soát của IRMS}\label{tab:appendix-component-rationale-adapters}\\
\toprule
\textbf{Thành phần / loại hình} & \textbf{Lý do dùng trong IRMS} & \textbf{Liên hệ trong kiến trúc} \\
\midrule
\endfirsthead
\toprule
\textbf{Thành phần / loại hình} & \textbf{Lý do dùng trong IRMS} & \textbf{Liên hệ trong kiến trúc} \\
\midrule
\endhead
\textsf{PaymentGateway} & Cô lập xử lý thanh toán sau một interface nội bộ, giúp \textsf{Billing \& Settlement Service} không phụ thuộc trực tiếp vào một cơ chế thanh toán cụ thể. & Là adapter biên của vùng thanh toán và xuất hiện trong kiến trúc tổng quan, C\&C View và sơ đồ lớp. \\
\addlinespace
\textsf{Notification}\-\textsf{Channel}\-\textsf{Adapter} & Tách lựa chọn kênh thông báo khỏi logic nghiệp vụ chính, cho phép thay đổi kênh gửi mà không sửa lõi thông báo. & Là adapter biên của vùng thông báo và liên kết với các luồng sự kiện nền. \\
\addlinespace
\textsf{Receipt}\-\textsf{Delivery}\-\textsf{Adapter} & Tách phát hành biên lai khỏi tuyến thanh toán lõi; biên lai giấy hoặc điện tử có thể xử lý sau khi trạng thái thanh toán đã rõ ràng. & Liên kết với \textsf{BillingReceiptService} và các luồng phát hành biên lai. \\
\addlinespace
Phân quyền theo vai trò và \textsf{AuthorizationPolicy} & Kiểm soát quyền theo vai trò và theo điều kiện nghiệp vụ, đặc biệt với hoàn tiền, thay đổi giá, điều chỉnh tồn kho và truy cập dữ liệu quản trị. & Gắn với yêu cầu bảo mật ở Chương~1 và xuất hiện trong kiến trúc tổng quan, C\&C View và thiết kế chi tiết. \\
\addlinespace
\textsf{AuditLog} & Ghi nhận thao tác nhạy cảm để hỗ trợ truy vết, đối soát và kiểm toán vận hành. & Dữ liệu kiểm toán là nhóm bảng logic trong cùng \textsf{PostgreSQL}; các activity diagram quản trị cần thể hiện điểm ghi nhận này. \\
\bottomrule
\end{longtable}
\normalsize
}

\subsection{Bảng thuật ngữ và ký hiệu viết tắt}\label{appendix:abbreviations}

Bảng này thống nhất cách hiểu các ký hiệu viết tắt xuất hiện trong tài liệu. Tên riêng hoặc thuật ngữ chuẩn được giữ nguyên; phần diễn giải dùng tiếng Việt để tránh nhầm lẫn khi đọc giữa các chương.

{\small
\setlength{\LTleft}{0pt}
\setlength{\LTright}{0pt}
\renewcommand{\arraystretch}{1.12}
\begin{longtable}{>{\raggedright\arraybackslash}p{2.8cm} >{\raggedright\arraybackslash}p{4.2cm} >{\raggedright\arraybackslash}p{6.8cm}}
\caption{Thuật ngữ và ký hiệu viết tắt dùng trong tài liệu}\label{tab:appendix-abbreviations}\\
\toprule
\textbf{Ký hiệu} & \textbf{Tên đầy đủ} & \textbf{Cách hiểu trong IRMS} \\
\midrule
\endfirsthead
\toprule
\textbf{Ký hiệu} & \textbf{Tên đầy đủ} & \textbf{Cách hiểu trong IRMS} \\
\midrule
\endhead
ADR & Architecture Decision Record & Hồ sơ quyết định kiến trúc, dùng để ghi lại lựa chọn quan trọng và hệ quả của lựa chọn đó. \\
API & Application Programming Interface & Hợp đồng giao tiếp giữa giao diện, gateway và các service. \\
C\&C & Component-and-Connector & Góc nhìn thành phần và kết nối, mô tả thành phần runtime, port, interface và connector. \\
IAM & Identity and Access Management & Vùng định danh và kiểm soát truy cập của hệ thống. \\
RBAC & Role-Based Access Control & Cơ chế phân quyền theo vai trò người dùng. \\
POS & Point of Sale & Điểm bán hàng hoặc giao diện thao tác bán hàng tại nhà hàng. \\
KDS & Kitchen Display System & Màn hình hiển thị bếp để nhận và cập nhật phiếu chế biến. \\
REST & Representational State Transfer & Kiểu giao tiếp đồng bộ dùng cho các thao tác cần phản hồi trực tiếp. \\
AMQP & Advanced Message Queuing Protocol & Giao thức hàng đợi dùng với \textsf{RabbitMQ}. \\
JDBC & Java Database Connectivity & Cơ chế kết nối từ ứng dụng Java đến \textsf{PostgreSQL}. \\
SOLID & Năm nguyên tắc thiết kế hướng đối tượng & Nhóm nguyên tắc giúp phân tách trách nhiệm, mở rộng có kiểm soát và đảo ngược phụ thuộc. \\
UI & User Interface & Giao diện người dùng. \\
UC & Use Case & Ca sử dụng được đặc tả trong Chương~2. \\
\bottomrule
\end{longtable}
\normalsize
}

\subsection{Phân công công việc}
Phần phân công dưới đây trình bày phạm vi đóng góp chính của từng thành viên. Các nội dung vẫn cần được rà soát chéo để bảo đảm mạch liên kết giữa yêu cầu, kiến trúc, thiết kế chi tiết và phần phụ lục.

\begin{table}[H]
\centering
\caption{Phân công công việc của Nhóm 8}
\small
\begin{tabularx}{\textwidth}{|p{1.9cm}|p{3.2cm}|X|}
\hline
\textbf{MSSV} & \textbf{Thành viên} & \textbf{Phạm vi đóng góp chính} \\
\hline
2311896 & Đỗ Quang Long &
Phụ trách tổng hợp bối cảnh nghiệp vụ, hệ thống hóa yêu cầu chức năng và phi chức năng, đồng thời phối hợp rà soát để bảo đảm các yêu cầu được phản ánh nhất quán trong các phần phân tích và thiết kế. \\
\hline
2311982 & Lê Hoàng Luân &
Phụ trách tổ chức cấu trúc tài liệu LaTeX, chuẩn hóa hình thức trình bày và tích hợp các thành phần nội dung, đồng thời phối hợp với các thành viên khác để bảo đảm bố cục báo cáo và hệ thống sơ đồ thống nhất. \\
\hline
2310990 & Lê Thúy Hiền &
Phụ trách xây dựng phần đặc tả use case và mô tả các luồng hoạt động nghiệp vụ chính, đồng thời phối hợp đối chiếu với phần yêu cầu và phần kiến trúc để giữ sự liên thông giữa mô tả nghiệp vụ và giải pháp thiết kế. \\
\hline
2213698 & Nguyễn Hải Trung &
Phụ trách các góc nhìn phân rã dịch vụ, thành phần và kết nối của hệ thống, đồng thời phối hợp với các phần khác để bảo đảm các quyết định kiến trúc bám sát yêu cầu nghiệp vụ và thống nhất với cấu trúc tổng thể của báo cáo. \\
\hline
2211384 & Tô Thế Hưng &
Phụ trách xây dựng các sơ đồ lớp và rà soát mô hình miền nghiệp vụ, đồng thời phối hợp đối chiếu với phần kiến trúc và phần hành vi để bảo đảm cấu trúc lớp phản ánh đúng các trách nhiệm và điểm mở rộng của hệ thống. \\
\hline
2310737 & Trần Nhật Đăng &
Phụ trách phần triển khai, phân bổ và phản tư kiến trúc, đồng thời phối hợp rà soát cách diễn đạt và mối liên kết giữa các chương để bảo đảm báo cáo có tính nhất quán và hoàn chỉnh ở mức toàn cục. \\
\hline
\end{tabularx}
\normalsize
\end{table}
